%! The Tangent Function — Unit 3, Lesson 3.8 — Guided Notes
\documentclass[10pt]{article}
\usepackage{precalculus-article}
\usepackage{precalculus-boxes}
\newcommand{\TallMath}[1]{$\displaystyle #1\rule[-1.4em]{0pt}{3.2em}$}

\begin{document}

\pageheader{Unit 3, Lesson 3.8}{Guided Notes}
\namedateperiod

\vspace{0.08in}

%----------------------------------------------------------------------
% Objectives
%----------------------------------------------------------------------
\begin{objectivebox}
\textbf{I will be able to\ldots}
\begin{enumerate}[leftmargin=1.5em, itemsep=4pt]
  \item \textbf{LO 3.8.A} --- Construct representations of $f(\theta)=\tan\theta$
    using the unit circle. \writeline
  \item \textbf{LO 3.8.B} --- Describe key characteristics of the tangent function
    (period, asymptotes, monotonicity, concavity). \writeline
  \item \textbf{LO 3.8.C} --- Describe additive and multiplicative transformations
    of the tangent function. \writeline
\end{enumerate}
\end{objectivebox}

\vspace{0.06in}

%----------------------------------------------------------------------
% Vocabulary
%----------------------------------------------------------------------
\begin{vocabbox}
\begin{description}[leftmargin=0pt, itemsep=6pt]
  \item[\termblanklong{tangent function}]
    The function $f(\theta) = $ \underline{\hspace{2cm}} that gives the
    \underline{\hspace{2.5cm}} of the terminal ray of angle $\theta$ on the unit circle.

  \item[\termblanklong{vertical asymptote}]
    A vertical line $\theta = a$ where the function values grow without
    \underline{\hspace{1.8cm}}.  For $\tan\theta$, these occur where
    $\cos\theta = $ \underline{\hspace{1cm}}.

  \item[\termblanklong{period (tangent)}]
    The tangent function has period \underline{\hspace{1cm}} because the
    slope of the terminal ray repeats every
    \underline{\hspace{2.2cm}}.
\end{description}
\end{vocabbox}

\vspace{0.06in}

%----------------------------------------------------------------------
% Hook
%----------------------------------------------------------------------
\begin{hookbox}
\textbf{Entry Question:} A lighthouse beam rotates at a constant rate.  The distance
along a straight shoreline to the point illuminated by the beam is a function of the
angle $\theta$ the beam makes with the perpendicular.

\vspace{4pt}
\textbf{Q1.} If the lighthouse is 1 unit from the shore and the angle is $\theta=\pi/4$,
how far along the shore does the beam shine?  (Hint: think about the tangent ratio in a
right triangle.)
\writelines{2}

\textbf{Q2.} What happens as $\theta$ approaches $\pi/2$?  Sketch your prediction on the
axes below.

\vspace{4pt}
\begin{center}
\begin{tikzpicture}[scale=0.75]
  \draw[->] (-0.3,0) -- (4.2,0) node[right]{\small $\theta$};
  \draw[->] (0,-0.3) -- (0,3.2) node[above]{\small distance};
  \draw[dashed, slate] (3.14,-.2) -- (3.14,3.2)
        node[above,font=\tiny]{$\pi/2$};
  \foreach \x/\l in {1.05/{$\pi/6$}, 2.09/{$\pi/3$}}
        \draw (\x,0.05)--(\x,-0.05) node[below,font=\tiny]{\l};
\end{tikzpicture}
\end{center}

\end{hookbox}

\vspace{0.06in}

%----------------------------------------------------------------------
% LO 3.8.A — Tangent from the Unit Circle
%----------------------------------------------------------------------
\begin{notesbox}{LO 3.8.A --- Constructing $\tan\theta$ from the Unit Circle}

\textbf{Definition:} For any angle $\theta$ in standard position, the terminal ray
passes through $(\cos\theta,\sin\theta)$.  The slope of that ray is:
\[
  \tan\theta = \frac{\blank{1.5cm}}{\blank{1.5cm}}
  \quad\text{provided } \cos\theta \neq \blank{0.8cm}
\]

\vspace{6pt}
\textbf{Complete the table} using $\sin$ and $\cos$ values from your warm-up.

\vspace{4pt}
\begin{center}
\renewcommand{\arraystretch}{1.55}
\begin{tabular}{|c|c|c|c|}
\hline
$\theta$ & $\sin\theta$ & $\cos\theta$ & $\tan\theta = \sin\theta/\cos\theta$ \\
\hline
$0$           & $0$             & $1$             & \blank{2.2cm} \\
\hline
$\pi/6$       & $1/2$           & $\sqrt{3}/2$    & \blank{2.2cm} \\
\hline
$\pi/4$       & $\sqrt{2}/2$    & $\sqrt{2}/2$    & \blank{2.2cm} \\
\hline
$\pi/3$       & $\sqrt{3}/2$    & $1/2$           & \blank{2.2cm} \\
\hline
$\pi/2$       & $1$             & $0$             & \blank{2.2cm} \\
\hline
$2\pi/3$      & $\sqrt{3}/2$    & $-1/2$          & \blank{2.2cm} \\
\hline
$3\pi/4$      & $\sqrt{2}/2$    & $-\sqrt{2}/2$   & \blank{2.2cm} \\
\hline
$\pi$         & $0$             & $-1$            & \blank{2.2cm} \\
\hline
\end{tabular}
\end{center}

\vspace{6pt}
\textbf{Annotate the pre-drawn graph} below with the values from your table.
Mark with an ``$\times$'' where vertical asymptotes occur.

\vspace{4pt}
\begin{center}
\begin{tikzpicture}[xscale=1.8, yscale=0.6]
  % axes
  \draw[->] (-1.8,0) -- (4.0,0) node[right]{\small $\theta$};
  \draw[->] (0,-3.3) -- (0,3.3) node[above]{\small $y$};
  % x-axis ticks
  \foreach \x/\l in {-1.571/$-\tfrac{\pi}{2}$, 0/$0$,
                      1.571/$\tfrac{\pi}{2}$, 3.142/$\pi$} {
    \draw (\x,0.08)--(\x,-0.08) node[below,font=\scriptsize]{\l};
  }
  % asymptotes
  \draw[dashed, slate] (-1.571,-3.2)--(-1.571,3.2);
  \draw[dashed, slate] ( 1.571,-3.2)--( 1.571,3.2);
  % pre-drawn tangent curves
  \draw[navy, thick, domain=-1.45:-0.0, samples=80, smooth]
        plot (\x, {tan(deg(\x))});
  \draw[navy, thick, domain=0.0:1.45, samples=80, smooth]
        plot (\x, {tan(deg(\x))});
  \draw[navy, thick, domain=1.70:3.10, samples=80, smooth]
        plot (\x, {tan(deg(\x))});
  % y-axis labels
  \foreach \y in {-2,-1,1,2}
    \draw (0.04,\y)--(-0.04,\y) node[left,font=\scriptsize]{$\y$};
\end{tikzpicture}
\end{center}

\end{notesbox}

\vspace{0.06in}

%----------------------------------------------------------------------
% LO 3.8.B — Key Characteristics
%----------------------------------------------------------------------
\begin{notesbox}{LO 3.8.B --- Key Characteristics of $y = \tan\theta$}

\renewcommand{\arraystretch}{1.6}
\begin{tabularx}{\linewidth}{@{} p{0.30\linewidth} X @{}}
\hline
\textbf{Characteristic} & \textbf{Description / Value} \\
\hline
Period & \blank{5cm} \\
Vertical asymptotes & $\theta = $ \blank{4cm} for integer $k$ \\
Zeros & $\theta = $ \blank{4cm} for integer $k$ \\
Monotonicity & The function is \blank{2.5cm} on every interval between consecutive asymptotes \\
Concavity change & Concave \blank{1.5cm} for $\theta<0$; concave \blank{1.5cm} for $\theta>0$ (within $(-\pi/2, \pi/2)$) \\
\hline
\end{tabularx}

\vspace{8pt}
\textbf{Why period $= \pi$:} After a half-revolution the terminal ray points in the
opposite direction, so both rise and run \underline{\hspace{4cm}}, giving the same
\underline{\hspace{2cm}}.

\end{notesbox}

\vspace{0.06in}

%----------------------------------------------------------------------
% LO 3.8.C — Transformations
%----------------------------------------------------------------------
\begin{notesbox}{LO 3.8.C --- Transformations of $y = \tan\theta$}

For $g(\theta) = a\cdot\tan(b(\theta + c)) + d$:

\vspace{4pt}
\renewcommand{\arraystretch}{1.6}
\begin{tabularx}{\linewidth}{@{} p{0.12\linewidth} p{0.32\linewidth} X @{}}
\hline
\textbf{Param.} & \textbf{Effect} & \textbf{Example} \\
\hline
$a$ & Vertical dilation by \blank{1cm} & $2\tan\theta$: outputs $\times 2$ \\
$b$ & Period becomes \blank{1.5cm} & $\tan(2\theta)$: period $= \pi/2$ \\
$c$ & Phase shift of \blank{1.5cm} & $\tan(\theta-\pi/4)$: shift right $\pi/4$ \\
$d$ & Vertical translation by \blank{1cm} & $\tan\theta + 3$: shift up $3$ \\
\hline
\end{tabularx}

\vspace{8pt}
\textbf{Finding new asymptotes:} Asymptotes occur when the \emph{inside argument}
equals $\pm\pi/2 + k\pi$.

\vspace{4pt}
\textbf{Example:} Find the asymptotes of $g(\theta) = \tan(2\theta)$.

Set $2\theta = \dfrac{\pi}{2} + k\pi$ \quad$\Rightarrow$\quad $\theta = $ \blank{3cm}

\end{notesbox}

\vspace{0.06in}

%----------------------------------------------------------------------
% Practice
%----------------------------------------------------------------------
\begin{practicebox}

\textbf{Guided Practice.}\quad For each function below, state the period and
the location of the vertical asymptote closest to $\theta = 0$ (with $\theta > 0$).

\vspace{4pt}
\renewcommand{\arraystretch}{1.6}
\begin{tabularx}{\linewidth}{@{} p{0.40\linewidth} X X @{}}
\hline
\textbf{Function} & \textbf{Period} & \textbf{Asymptote ($\theta > 0$)} \\
\hline
$g(\theta) = \tan(3\theta)$ & \blank{2.2cm} & \blank{2.4cm} \\
$g(\theta) = \tan\!\left(\tfrac{1}{2}\theta\right)$ & \blank{2.2cm} & \blank{2.4cm} \\
$g(\theta) = 4\tan(\theta) + 1$ & \blank{2.2cm} & \blank{2.4cm} \\
\hline
\end{tabularx}

\end{practicebox}

\end{document}
